% GENERATED FILE. Edit canonical lessons, then run npm run generate:books.
% canonical-source-hash: fnv1a64:c3000e2d6e84f16e
% canonical-lessons: BN-C68-mati, BN-C68-batas, BN-C68-akash, BN-C68-kukur, BN-C68-biral

\chapter{Soil, Wind, Sky, Dog, Cat}
\label{ch:bn-soil-wind-sky-dog-cat}
\hlchaptermodality{68}

\begin{chapteropening}
\textbf{By the end of this chapter:} \emph{I can say soil, earth, wind, the sky, a dog and a cat.}

Complete the last lesson of chapter 68: I can say soil, earth, wind, the sky, a dog and a cat.
\end{chapteropening}

\section[\texorpdfstring{\bn{মাটি} (māṭi)}{মাটি (māṭi)}]{\bn{মাটি} (māṭi) — soil, earth}

\label{lesson:BN-C68-mati}



\begin{quote}
Before the new one: say the Bengali for fire, then the Bengali for ice.
\end{quote}

\subsection*{You'll want to know: \bn{মাটি}}
\textbf{\bn{মাটি}} — \emph{māṭi} — \textquotedblleft{}soil, earth\textquotedblright{}.

\begin{cousinweb}
\textbf{\bn{মাটি}} (\emph{māṭi}) is a tadbhava, a word inherited from Sanskrit and reshaped by sound change. It comes from \emph{mṛttikā}, \textquotedblleft{}earth, clay\textquotedblright{} (CDIAL 10286).
\end{cousinweb}

\begin{grammarlens}[title={What you've built}]
One more word for this chapter.
\end{grammarlens}

\subsection*{Your turn}
\begin{itemize}
\raggedright
  \item \emph{Say it:} \emph{māṭi}
  \item \emph{Say it:} \emph{māṭi}, once more
  \item \emph{Say it:} say \emph{māṭi}
  \item \emph{Recall:} say the Bengali for fire, then the Bengali for ice, then say \emph{māṭi} again
\end{itemize}

\begin{tcolorbox}[breakable,colback=teal!4,colframe=teal!35!black,title={Before you move on}]
What does \bn{মাটি} mean? (\textquotedblleft{}Soil, earth\textquotedblright{}.) Say it once more.
\end{tcolorbox}
\section[\texorpdfstring{\bn{বাতাস} (bātās)}{বাতাস (bātās)}]{\bn{বাতাস} (bātās) — wind}

\label{lesson:BN-C68-batas}



\begin{quote}
Before the new one: say the Bengali for ice, then the Bengali for soil.
\end{quote}

\subsection*{You'll want to know: \bn{বাতাস}}
\textbf{\bn{বাতাস}} — \emph{bātās} — \textquotedblleft{}wind\textquotedblright{}.

\begin{cousinweb}
\textbf{\bn{বাতাস}} (\emph{bātās}) is a tadbhava, a word inherited and reshaped by sound change. It comes from \emph{*vātatrāsa}, \textquotedblleft{}wind-storm\textquotedblright{}, a form CDIAL reconstructs from its descendants (CDIAL 11493). It is a compound of \emph{vāta}, \textquotedblleft{}wind\textquotedblright{}, and \emph{trāsa}, \textquotedblleft{}fear\textquotedblright{}.
\end{cousinweb}

\begin{grammarlens}[title={What you've built}]
One more word for this chapter.
\end{grammarlens}

\subsection*{Your turn}
\begin{itemize}
\raggedright
  \item \emph{Say it:} \emph{bātās}
  \item \emph{Say it:} \emph{bātās}, once more
  \item \emph{Say it:} say \emph{bātās}
  \item \emph{Recall:} say the Bengali for ice, then the Bengali for soil, then say \emph{bātās} again
\end{itemize}

\begin{tcolorbox}[breakable,colback=teal!4,colframe=teal!35!black,title={Before you move on}]
What does \bn{বাতাস} mean? (\textquotedblleft{}Wind\textquotedblright{}.) Say it once more.
\end{tcolorbox}
\section[\texorpdfstring{\bn{আকাশ} (ākāsh)}{আকাশ (ākāsh)}]{\bn{আকাশ} (ākāsh) — the sky}

\label{lesson:BN-C68-akash}



\begin{quote}
Before the new one: say the Bengali for soil, then the Bengali for wind.
\end{quote}

\subsection*{You'll want to know: \bn{আকাশ}}
\textbf{\bn{আকাশ}} — \emph{ākāsh} — \textquotedblleft{}the sky\textquotedblright{}.

\begin{cousinweb}
\textbf{\bn{আকাশ}} (\emph{ākāsh}) is Sanskrit \emph{ākāśa}, \textquotedblleft{}sky\textquotedblright{}, taken over whole as a learned loan, a tatsama.
\end{cousinweb}

\begin{grammarlens}[title={What you've built}]
One more word for this chapter.
\end{grammarlens}

\subsection*{Your turn}
\begin{itemize}
\raggedright
  \item \emph{Say it:} \emph{ākāsh}
  \item \emph{Say it:} \emph{ākāsh}, once more
  \item \emph{Say it:} say \emph{ākāsh}
  \item \emph{Recall:} say the Bengali for soil, then the Bengali for wind, then say \emph{ākāsh} again
\end{itemize}

\begin{tcolorbox}[breakable,colback=teal!4,colframe=teal!35!black,title={Before you move on}]
What does \bn{আকাশ} mean? (\textquotedblleft{}The sky\textquotedblright{}.) Say it once more.
\end{tcolorbox}
\section[\texorpdfstring{\bn{কুকুর} (kukur)}{কুকুর (kukur)}]{\bn{কুকুর} (kukur) — a dog}

\label{lesson:BN-C68-kukur}



\begin{quote}
Before the new one: say the Bengali for wind, then the Bengali for the sky.
\end{quote}

\subsection*{You'll want to know: \bn{কুকুর}}
\textbf{\bn{কুকুর}} — \emph{kukur} — \textquotedblleft{}a dog\textquotedblright{}.

\begin{cousinweb}
\textbf{\bn{কুকুর}} (\emph{kukur}) is a tadbhava, a word inherited from Sanskrit and reshaped by sound change. It comes from \emph{kukkura}, \textquotedblleft{}dog\textquotedblright{} (CDIAL 3329). CDIAL calls the word imitative of the animal's sound.
\end{cousinweb}

\begin{grammarlens}[title={What you've built}]
One more word for this chapter.
\end{grammarlens}

\subsection*{Your turn}
\begin{itemize}
\raggedright
  \item \emph{Say it:} \emph{kukur}
  \item \emph{Say it:} \emph{kukur}, once more
  \item \emph{Say it:} say \emph{kukur}
  \item \emph{Recall:} say the Bengali for wind, then the Bengali for the sky, then say \emph{kukur} again
\end{itemize}

\begin{tcolorbox}[breakable,colback=teal!4,colframe=teal!35!black,title={Before you move on}]
What does \bn{কুকুর} mean? (\textquotedblleft{}A dog\textquotedblright{}.) Say it once more.
\end{tcolorbox}
\section[\texorpdfstring{\bn{বিড়াল} (biṛāl)}{বিড়াল (biṛāl)}]{\bn{বিড়াল} (biṛāl) — a cat}

\label{lesson:BN-C68-biral}



\begin{quote}
Before the new one: say the Bengali for the sky, then the Bengali for a dog.
\end{quote}

\subsection*{You'll want to know: \bn{বিড়াল}}
\textbf{\bn{বিড়াল}} — \emph{biṛāl} — \textquotedblleft{}a cat\textquotedblright{}.

\begin{cousinweb}
\textbf{\bn{বিড়াল}} (\emph{biṛāl}) is a tadbhava, a word inherited from Sanskrit and reshaped by sound change. It comes from \emph{biḍāla}, \textquotedblleft{}cat\textquotedblright{} (CDIAL 9237). CDIAL writes the Bengali form \emph{birāl}.
\end{cousinweb}

\begin{grammarlens}[title={What you've built}]
One more word for this chapter.
\end{grammarlens}

\subsection*{Your turn}
\begin{itemize}
\raggedright
  \item \emph{Say it:} \emph{biṛāl}
  \item \emph{Say it:} \emph{biṛāl}, once more
  \item \emph{Say it:} say \emph{biṛāl}
  \item \emph{Recall:} say the Bengali for the sky, then the Bengali for a dog, then say \emph{biṛāl} again
\end{itemize}

\begin{tcolorbox}[breakable,colback=teal!4,colframe=teal!35!black,title={Before you move on}]
What does \bn{বিড়াল} mean? (\textquotedblleft{}A cat\textquotedblright{}.) Say it once more.
\end{tcolorbox}
